% Exact uniform j=3,4,5 residual certificates, independently verified by
% /root/pencil. Dependencies: finding/proof, amsmath/amssymb, longtable,
% booktabs, \BR. P-01, P-06, R-01 and R-04 must appear elsewhere.

\subsection{How the finite certificates cover the Grassmannian}

For this subsection \(n=4,b=2\) and \(U\) is either polynomial
residual image. Use the ordered list
\[
 \mathcal P=((1,2),(1,3),(1,4),(2,3),(2,4),(3,4))
\]
for both row and column pairs. Label the rectangle
\((\mathcal P_a,\mathcal P_b)\) by \(6a+b\), with
\(0\leq a,b\leq5\). Thus certificate subset labels range from
zero through \(35\); matrix-entry indices remain \(1,\ldots,4\).

For a coordinate subset \(S\) of cardinality \(j\), let
\(E_S\) be the span of its used matrix variables and put
\(m= \dim E_S\). Every slice is a nonsingular signed matching
on its four-coordinate support, so the last-factor flattening
has rank \(m\). For the residual tensor
\(T_S\in A_S\otimes E_S\otimes E_S\), \(\dim A_S=j\), use
\[
 K_{S,p}:\bigwedge^p A_S\otimes E_S^*
 \longrightarrow \bigwedge^{p+1}A_S\otimes E_S.
\]
Its pure-tensor cost is \(c_{j,p}=\binom{j-1}{p}\).
An integer minor of order \(d\) and nonzero determinant
therefore proves
\[
 \BR(T_S)\geq \left\lceil d/c_{j,p}\right\rceil
\]
over every field in which that determinant remains nonzero.
The proof of the cost is direct: for a pure tensor
\(a\otimes b\otimes c\), the map is the product of evaluation
by \(b\), the wedge map \(a\wedge-\), and the fixed output
vector \(c\). The wedge map has rank \(\binom{j-1}{p}\).
Rank subadditivity and closed determinantal loci prove the
border-rank bound.

Row permutations, column permutations, and transposition form
a group of order \(24^2\cdot2=1152\) preserving border rank
of coordinate residual tensors. In the determinant case they
may change individual slice signs, which are absorbed by
an invertible diagonal map on \(A_S\).
The complete orbit counts are
\[
 \begin{array}{c|r|r}
 j&\binom{36}{j}&\text{coordinate orbits}\\ \hline
 3&7140&21\\
 4&58905&99\\
 5&376992&454.
 \end{array}
\]
The source program \path{coordinate_koszul.py} constructs
all group images and all orbit representatives. Its modular
elimination is used only to locate candidate minors.
The independent certificate package
\path{global_j345_certificate.json} records the actual
row and column indices and integer determinant of every
minor needed below.

The independent, standard-library-only checker
\path{verify_global_bounds.py} rebuilds every minor
from the defining wedge map and verifies its determinant by
fraction-free integer elimination. It also enumerates
\emph{every} \(j\)-subset of the \(36\) rectangles and
checks that every subset whose ordinary support is below
the claimed threshold lies in one of the listed low-support
orbits. Thus no generic-rank sampling or unverified orbit
representative assumption is used. Finally P-01 transfers
the coordinate lower bound to all planes in the Grassmannian.

\begin{finding}{R-02}{the three-dimensional polynomial residual minimum}
For both polynomial families over every algebraically closed
field,
\[
 \mu_3(U_{4,2})=8.
\]
Consequently R-01 gives
\(\BR(\mathcal D_4),\BR(\mathcal P_4)\geq41\).
\end{finding}
\begin{proof}
Every three-rectangle subset with at least eight variables
is settled by its ordinary flattening. All remaining subsets
form a single orbit, represented by \(S=\{0,1,3\}\); it
contains \(48\) coordinate subsets and uses six variables.
For \(p=1\), where the pure cost is two, the determinant
tensor has an order-\(18\) minor of determinant \(-4\);
the permanent tensor has an order-\(16\) minor of determinant
\(1\). Hence the required lower bound eight holds for the
permanent in every characteristic and for the determinant
away from characteristic two. In characteristic two the
two tensors coincide, so the permanent certificate also
proves the determinant assertion.

For the upper bound use the anchored shape with one row
vertex, three column vertices, and its three edges.
P-06 gives border rank \(3+1+3+1=8\). This shape occurs
at \(S=\{0,1,2\}\). Hence the minimum is exactly eight.
Substitution gives \(36-3+8=41\).
\end{proof}

\begin{finding}{R-03}{the four-dimensional polynomial residual minimum}
For both polynomial families over every algebraically closed
field,
\[
 \mu_4(U_{4,2})=9.
\]
The resulting substitution bound is again \(41\).
\end{finding}
\begin{proof}
Subsets with at least nine variables are settled by the
ordinary flattening. Exactly \(900\) coordinate subsets
remain; their four orbits have the following \(p=1\)
certificates. The pure cost is three; each pair in the
table is (minor order, integer determinant).
\[
 \begin{array}{c|r r|r r}
 S&\multicolumn{2}{c|}{\text{determinant}}&
    \multicolumn{2}{c}{\text{permanent}}\\ \hline
 \{0,1,2,3\}&30&-4&30&1\\
 \{0,1,3,6\}&29&4&28&-1\\
 \{0,1,4,5\}&32&-1&32&-1\\
 \{0,1,9,21\}&31&1&31&-1
 \end{array}
\]
All permanent minors are units; the determinant minors
are nonzero away from characteristic two. In characteristic
two use equality of the two coefficient tensors. These
bounds exceed nine on the exceptional orbits and establish
the uniform lower bound.
The complete anchored \(2\times2\) shape
\(S=\{0,1,6,7\}\) has border rank \(4+2+2+1=9\) by P-06.
Thus the minimum is nine, and substitution gives
\(36-4+9=41\).
\end{proof}

\begin{finding}{R-05}{the five-dimensional minimum and the bound forty-two}
For both polynomial families over every algebraically closed
field, in the integral coefficient-tensor convention,
\[
 \boxed{\mu_5(U_{4,2})=11},\qquad
 \boxed{\BR(\mathcal D_4)\geq42,\quad
        \BR(\mathcal P_4)\geq42}.
\]
Over characteristic zero these are ordinary CP border-rank
lower bounds for the usual determinant and permanent
polynomial tensors, and therefore also symmetric
border-rank lower bounds.
\end{finding}
\begin{proof}
Ordinary flattenings settle every coordinate subset with
at least eleven variables. The remaining subsets have
exactly \(26\) orbits. Table~\ref{tab:five-minors}
gives one sufficient certificate for each family on every
one of these orbits. Its entries are triples
\((p;d;\Delta)\), where \(p\) is the Koszul exterior degree,
\(d\) the minor order, and \(\Delta\) its exact integer
determinant. For \(p=1\) the pure cost is four, and every
listed \(d\geq41\) certifies at least eleven. For \(p=2\)
the pure cost is six, and every listed \(d\geq61\)
does the same.

All listed permanent determinants are \(1\) or \(-1\).
The single permanent exception \(\{0,1,3,6,7\}\) is
precisely \(T_5\) of P-08; the unit-minor prolongation
obstruction R-04 gives its lower bound eleven in every
characteristic. All determinant minors in the table are
units times powers of two. They therefore prove the
determinant bound in every characteristic other than two,
and in characteristic two the permanent certificates apply
because the two integral tensors coincide. The
order-\(61\), determinant-\(-1\) certificate in the first
row is included specifically to remove the characteristic
three exceptional prime of a larger, unnecessary minor.

This proves the uniform lower bound. An anchored five-edge
shape in a \(2\times3\) grid, for example
\(S=\{0,1,2,6,7\}\), has border rank
\(5+2+3+1=11\) by P-06. Hence the minimum is attained and
equals eleven for both signs. R-01 now gives
\(36-5+11=42\).
\end{proof}

{\small
\begin{longtable}{@{}lcc@{}}
\caption{Complete low-support five-plane certificates.
Entries are \((p;d;\Delta)\).}
\label{tab:five-minors}\\
\toprule
Rectangle subset \(S\)&Determinant&Permanent\\
\midrule
\endfirsthead
\toprule
Rectangle subset \(S\)&Determinant&Permanent\\
\midrule
\endhead
\(\{0,1,2,3,4\}\) & \(2;61;-1\) & \(2;72;1\) \\
\(\{0,1,2,3,6\}\) & \(1;44;4\) & \(1;44;-1\) \\
\(\{0,1,2,3,8\}\) & \(1;46;-4\) & \(1;46;-1\) \\
\(\{0,1,2,3,9\}\) & \(1;46;2\) & \(1;46;-1\) \\
\(\{0,1,2,3,10\}\) & \(1;47;4\) & \(1;47;1\) \\
\(\{0,1,2,6,18\}\) & \(1;42;-2\) & \(1;41;1\) \\
\(\{0,1,2,9,21\}\) & \(1;48;1\) & \(1;48;-1\) \\
\(\{0,1,3,6,7\}\) & \(1;41;2\) & \(\text{R-04}\) \\
\(\{0,1,3,6,8\}\) & \(1;46;4\) & \(1;45;1\) \\
\(\{0,1,3,6,13\}\) & \(1;43;-2\) & \(1;42;-1\) \\
\(\{0,1,3,6,18\}\) & \(2;71;64\) & \(2;66;1\) \\
\(\{0,1,3,6,19\}\) & \(1;42;-2\) & \(1;42;-1\) \\
\(\{0,1,3,6,20\}\) & \(1;47;-4\) & \(1;47;-1\) \\
\(\{0,1,3,6,24\}\) & \(1;44;4\) & \(1;44;1\) \\
\(\{0,1,3,6,25\}\) & \(1;46;2\) & \(1;46;-1\) \\
\(\{0,1,3,6,30\}\) & \(1;48;-4\) & \(1;47;1\) \\
\(\{0,1,3,8,10\}\) & \(1;47;4\) & \(1;46;1\) \\
\(\{0,1,3,8,20\}\) & \(1;47;-4\) & \(1;47;-1\) \\
\(\{0,1,4,5,6\}\) & \(1;48;-1\) & \(1;48;-1\) \\
\(\{0,1,4,5,8\}\) & \(1;46;-1\) & \(1;46;-1\) \\
\(\{0,1,4,8,20\}\) & \(1;46;1\) & \(1;46;1\) \\
\(\{0,1,4,11,23\}\) & \(1;50;-1\) & \(1;50;1\) \\
\(\{0,1,6,7,21\}\) & \(1;44;1\) & \(1;44;-1\) \\
\(\{0,1,6,9,19\}\) & \(1;43;1\) & \(1;43;1\) \\
\(\{0,1,6,15,27\}\) & \(1;48;-1\) & \(1;48;1\) \\
\(\{0,1,9,22,23\}\) & \(1;50;1\) & \(1;50;1\) \\
\bottomrule
\end{longtable}
}

The table is not a claim of optimal Koszul rank for every
orbit; it records enough exact minors for the uniform
bound. R-04 is needed because the direct Koszul certificates
alone leave the exceptional permanent five-plane at ten.
No priority assertion is made by R-02, R-03 or R-05;
the literature comparison is a separate part of the paper.

\paragraph{Algebraic-geometric form.}
The Grassmannian minima are avoided closed relative Segre
secant loci for the tautological residual tensor.
The Koszul maps are morphisms of locally free sheaves
formed with exterior powers, and their integer minors
define determinantal ideals. The units in the certificate
give open rank conditions after every base change.
The only residual characteristic issue is removed by the
identity of the two signs in characteristic two.
R-04 adds a proper Grassmannian incidence with the
first-prolongation determinantal loci, the low-degree
Hilbert-function test of border apolarity.
